\section{SF-E837. Every formal square-thickening point modification is
excluded}\label{SF-E837:sf-e837.-every-formal-square-thickening-point-modification-is-excluded}

Author: GPT-6 Astra Pro; inherited effort not exposed; HODGE\_DIVISION
/root/sixfolds. Date: 2026-11-21.

\textbf{Status:} complete new proof checked in full by the author,
/root and /root/g34\_eightfolds. The root also read the complete
canonical conversion, with all seventeen display blocks unchanged.
It uses the opposite-order companion equations proved in E836
sections 2--3, whose separate whole-proof reviews have passed. It does not use E836's commuting-triple hypothesis. The
integrable scalar-cross-minor boundary remains valid; the new argument
uses a strictly stronger actual matrix constraint.

\subsection{1. Statement}\label{SF-E837:statement}

Let \(R=\mathbb C[[z_1,z_2,z_3,z_4]]\), and let \(\pi_+,\pi_-\) be
constant rank-two bivectors with complementary supports. Let
\(u_1,u_2,u_3,u_4\) be any formal regular parameters, possibly nonlinear
and mixing both planes. Set

\[
 I=(u_1^2,u_2^2,u_3^2,u_4^2),
 \qquad A=R/I,
 \qquad\mathfrak n=(\bar u_1,\bar u_2,\bar u_3,\bar u_4)\subset A.
 \tag{E837.1}
\]

Consider a finite module with an exact sequence

\[
 0\longrightarrow A^n\longrightarrow M\longrightarrow R^r
     \longrightarrow A^m\longrightarrow0,
 \qquad m+n>0.
 \tag{E837.2}
\]

Then \(M\) has no pair of flat first-order module deformations for the
two multiplication classes \(\pi_+,\pi_-\). The assertion holds for
every positive multiplicity and every formal orientation of the square
thickening. It consequently excludes weak balanced semiregularity for a
coherent sheaf with such a nontrivial local point modification and any
specified coherent effective finite translation/tensor action in E819's
scope.

The proof has two ingredients: a universal matrix lemma over the
length-sixteen square algebra and the constant terms of the two
transformed bivectors. It does not identify a vanishing coefficient
cross minor with a vanishing Jacobian.

\subsection{2. Four monomial annihilators force a square matrix to be
invertible}\label{SF-E837:four-monomial-annihilators-force-a-square-matrix-to-be-invertible}

We write \(t_i\) for \(\bar u_i\), identifying
\(A=\mathbb C[t_1,t_2,t_3,t_4]/(t_1^2,t_2^2,t_3^2,t_4^2)\). Its
degree-one and degree-two pieces are

\[
 V=\operatorname{span}_{\mathbb C}(t_1,t_2,t_3,t_4),
 \qquad W=\operatorname{span}_{\mathbb C}(t_it_j: i<j),
 \qquad \dim V=4,\quad\dim W=6.
 \tag{E837.3}
\]

\textbf{Matrix lemma.} Let \(C\in\operatorname{Mat}_m(A)\). If the
cokernel of \(C:A^m\to A^m\) is annihilated by at least four distinct
monomials \(t_it_j\), then \(C\) is invertible.

First split off all unit pivots by invertible row and column operations
over the local ring. If a unit entry exists, scale it to one, clear its
row and column, and repeat on the remaining square block. The process
terminates in a block decomposition with an identity block and a square
block \(C'\) of size \(k\), all of whose entries belong to
\(\mathfrak n\). The cokernel is unchanged up to isomorphism. If
\(k=0\), the claim holds. Assume \(k>0\); the four monomials still
annihilate \(\operatorname{coker}C'\).

Work modulo \(\mathfrak n^3\) and write the matrix as \(C_1+C_2\), where
the entries are homogeneous of degrees one and two. Let

\[
 L:\mathbb C^k\longrightarrow V\otimes\mathbb C^k,
 \qquad L(v)=C_1v,
 \qquad a=\operatorname{rank}_{\mathbb C}L\le k.
 \tag{E837.4}
\]

Any input modulo \(\mathfrak n^2\) is a sum of a constant vector \(v_0\)
and a linear vector \(v_1\). Its image modulo \(\mathfrak n^3\) is
\(C_1v_0+C_2v_0+C_1v_1\). Such an image has no degree-one term precisely
when \(v_0\in\ker L\). Therefore the intersection of the image with the
pure degree-two target is exactly

\[
 \operatorname{im}(C'\bmod\mathfrak n^3)
       \cap(W\otimes\mathbb C^k)
     = C_2(\ker L)+ V\cdot\operatorname{im}L.
 \tag{E837.5}
\]

The second summand has dimension at most \(4a\), because it is the image
of \(V\otimes\operatorname{im}L\). The first has dimension at most
\(k-a\). Thus the dimension in (E837.5) is at most

\[
 k+3a\le4k.
 \tag{E837.6}
\]

Choose four distinct pair monomials that annihilate the cokernel and let
\(J\subset W\) be their four-dimensional span. Annihilation means that
each of these monomials times every constant target vector belongs to
the image of \(C'\). Hence \(J\otimes\mathbb C^k\), of dimension \(4k\),
is contained in (E837.5). Equality must hold throughout (E837.6):
\(a=k\), \(\ker L=0\), and

\[
 V\cdot\operatorname{im}L=J\otimes\mathbb C^k.
 \tag{E837.7}
\]

Write \(C_1=\sum_i t_i A_i\), with constant \(k\)-by-\(k\) matrices
\(A_i\). For any pair \(\{i,j\}\) whose monomial is not in the chosen
four, multiply \(C_1v\) by \(t_i\). Its coefficient at \(t_it_j\) is
\(A_jv\), which must be zero by (E837.7). Multiplication by \(t_j\)
similarly gives \(A_i=0\). There are exactly two omitted distinct edges
among the six pairs of four vertices. Their endpoints cover at least
three vertices. Consequently at most one of the matrices \(A_i\) can be
nonzero.

If the one possible index is \(s\), then \(C_1=t_sA_s\), and the
dimension of \(V\cdot\operatorname{im}L\) is at most
\(3\operatorname{rank}A_s\le3k\), since \(t_s^2=0\). This contradicts
(E837.7). Thus \(k>0\) was impossible, proving the matrix lemma. The
argument applies in every size; no simultaneous diagonalization or
matrix commutativity was assumed.

\subsection{3. The actual companion equations in the square
algebra}\label{SF-E837:the-actual-companion-equations-in-the-square-algebra}

If a deformation exists for either nonzero bivector, the E819
finite-cone trace gives \(16n=16m\). We may therefore assume \(n=m>0\).
For the regular generators \(f_i=u_i^2\), the actual module construction
and both matrix-order calculations of E836 give fixed matrices
\(B_{ij}\) and two matrices \(C_+,C_-\) in \(\operatorname{Mat}_m(A)\)
with

\[
 C_\pm B_{ij}=B_{ij}C_\pm=P^\pm_{ij}I_m,
 \qquad
 P^\pm_{ij}=4t_it_j\,a^\pm_{ij}(t),
 \tag{E837.8}
\]

where \(a^\pm_{ij}\) is the image of the actual formal function
\(\pi_\pm(du_i,du_j)\) in \(A\). These functions need not be constant.
Their values at the origin define alternating forms

\[
 p_{ij}=a^+_{ij}(0),\qquad q_{ij}=a^-_{ij}(0).
 \tag{E837.9}
\]

The derivative matrix of the formal parameter change is invertible.
Hence \(p\) and \(q\) are decomposable rank-two forms with complementary
supports, so

\[
 p\wedge p=q\wedge q=0,
 \qquad p\wedge q\ne0.
 \tag{E837.10}
\]

The scalar cross-minor consequences of the two matrix orders are
\(P^+_{ij}P^-_{kl}=P^-_{ij}P^+_{kl}\). For disjoint pairs, all higher
terms of the functions \(a^\pm\) disappear after multiplying by the top
squarefree monomial. That monomial is nonzero in \(A\). Thus

\[
 p_{ij}q_{kl}=q_{ij}p_{kl}
       \quad\text{whenever }\{i,j,k,l\}=\{1,2,3,4\}.
 \tag{E837.11}
\]

This reduction is legitimate even for a nonlinear formal parameter
change; it uses the explicit square-algebra socle, not a linearity
assumption on that change.

\subsection{4. The small-support case already contradicts scalar
minors}\label{SF-E837:the-small-support-case-already-contradicts-scalar-minors}

Let \(\mathcal E\) be the set of unordered pairs \(\{i,j\}\) for which
at least one of \(p_{ij},q_{ij}\) is nonzero. Suppose first that
\(|\mathcal E|\le3\). Among at most three edges there is at most one
pair of disjoint complementary edges: two distinct complementary pairs
require four different edges. If there is no complementary pair, the
mixed wedge is zero, contradicting (E837.10).

For the possible unique complementary pair, write the two coefficients
of \(p\) as \(a,b\) and those of \(q\) as \(c,d\), with a fixed
ordering. The decomposable identities (E837.10) give \(ab=cd=0\), since
this is the only possible term in each Pfaffian. Equation (E837.11)
gives \(ad=cb\). If their common value is \(z\), then

\[
 z^2=(ad)(cb)=(ab)(cd)=0.
 \tag{E837.12}
\]

These are complex numbers, so \(z=0\). The only possible mixed-wedge
coefficient is, up to sign, \(ad+cb=0\), again a contradiction.
Therefore

\[
 |\mathcal E|\ge4.
 \tag{E837.13}
\]

No analogous step is made for arbitrary coefficients in a nonreduced
ring; only the constant complex coefficients are used here.

\subsection{5. The large-support case forces an impossible scalar
ratio}\label{SF-E837:the-large-support-case-forces-an-impossible-scalar-ratio}

Choose a nonzero complex number \(\lambda\) outside the finite set of
roots of \(p_{ij}+\lambda q_{ij}\) for \(\{i,j\}\in\mathcal E\). Define
the actual matrix pencil

\[
 C_\lambda=C_++\lambda C_-,
 \qquad P^\lambda_{ij}=P^+_{ij}+\lambda P^-_{ij}.
 \tag{E837.14}
\]

Linearity of (E837.8) gives
\(C_\lambda B_{ij}=B_{ij}C_\lambda=P^\lambda_{ij}I_m\). This is an
algebraic matrix identity; the proof does not need to assume an
additional deformation for every parameter value.

For every edge in \(\mathcal E\),
\(P^\lambda_{ij}=4t_it_j(a^+_{ij}+\lambda a^-_{ij})\), and the factor in
parentheses is a unit of the local Artinian ring by the choice of
\(\lambda\). Consequently each \(t_it_j\) in at least four distinct
pairs annihilates \(\operatorname{coker}C_\lambda\). The matrix lemma
proves that \(C_\lambda\) is invertible.

Put \(D=C_+C_\lambda^{-1}\). Solving the first pencil equation for
\(B_{ij}\), and substituting in its plus equation, gives

\[
 D P^\lambda_{ij}=P^+_{ij}I_m.
 \tag{E837.15}
\]

Let \(D_0\) be the constant matrix modulo \(\mathfrak n\). The
degree-two component of (E837.15), for every edge in \(\mathcal E\), is

\[
 (p_{ij}+\lambda q_{ij})D_0=p_{ij}I_m.
 \tag{E837.16}
\]

One such nonzero coefficient forces \(D_0=d I_m\) for a complex scalar
\(d\). All remaining equations imply \(p=d(p+\lambda q)\) as complete
alternating forms, with zero coefficients off \(\mathcal E\) as well.
Since \(p\ne0\), \(d\ne0\). Since \(\lambda\ne0\), this makes \(p\) and
\(q\) linearly dependent, contradicting their complementary supports in
(E837.10). This completes the local theorem.

Both parts of the proof are necessary: the dense-support argument uses
the actual square matrices and their cokernels, while the small-support
argument uses only constant Pluecker and cross-minor identities. The
unrestricted scalar-cross-minor implication refuted by the integrable
boundary example is never invoked.

\subsection{6. Global consequence and actual excluded
objects}\label{SF-E837:global-consequence-and-actual-excluded-objects}

The global passage is precisely the reviewed E819 one. A balanced weakly
semiregular sheaf, with its specified coherent finite translation/tensor
action, makes the invariant evaluation vanish on the two balanced
cap-kernel classes. Forgetting the action and restricting to the
completed local ring gives flat first-order deformations for the two
constant bivectors. A rational algebraic \(B\) shear preserves those
bivector components. Thus a single nontrivial point modification of type
(E837.2) is impossible. The chosen formal parameters and individual
defect bases do not need equivariant structures.

There are actual ordinarily gluable balanced coherent sheaves in this
family before semiregularity is imposed. Take the rank-32 bundle \(H\)
of E815/E820, with full character \(32v_{1/2}\), and any finite
subscheme \(Z\) with the completed ideal (E837.1). It is an actual
algebraic subscheme: the ideal contains a power of the maximal ideal,
and finite quotients of the local ring and its completion agree. Choose
a quotient \(H\twoheadrightarrow\mathcal O_Z^m\), with \(1\le m\le32\),
let \(K\) be its kernel, and choose any extension

\[
 0\longrightarrow\mathcal O_Z^m\longrightarrow E\longrightarrow K
       \longrightarrow0.
 \tag{E837.17}
\]

As in E836, the extension space has dimension \(96m^2\), because the
quotient has length sixteen and a four-generator Koszul resolution. The
torsion of \(E\) is the displayed point module, its double dual is
\(H\), and its hull cokernel is the same point module. Therefore its
entire character is \(32v_{1/2}\). It is ordinarily gluable because it
is a coherent sheaf. Every such object is excluded from weak
semiregularity by the theorem, without any assertion that a prescribed
finite action actually exists.

In particular the parent's checked integrable boundary ideal, defined by
four squared mixed rational linear forms and having all scalar cross
minors zero but a nonzero mixed wedge, is now excluded by the stronger
matrix argument. Its counterexample to the scalar inference is
unchanged. This is a strict new family beyond the E836 commuting-triple
criterion, not a reclassification of that boundary example as a
deformation.

Arbitrary Artinian defect modules, common complete-intersection algebras
not isomorphic to this formal square thickening and not covered by E836,
unstable torsion-free objects, and unrelated perfect complexes remain
outside the theorem. The new proof uses no additional external analytic
or algebraic classification theorem. Its finite supplement only checks
the elementary edge combinatorics and selected exact coefficient
identities; the all-size filtered-matrix argument is the proof in
section 2.
